\section{Department of Corrections Profile}

The Department of Corrections (DOC) profile is the second validation vertical.
It demonstrates that the universal engine can evaluate a public custodial
institution without importing consumer-credit metrics.

\subsection{Profile Scope}

The DOC profile may cover a state department, facility, division, or program,
but the subject and period must be explicit. Facility-level evidence must not
be generalized to the entire department without a documented aggregation rule.

\subsection{Candidate Metric Families}

\begin{center}
\begin{tabular}{p{3.5cm}p{10cm}}
\toprule
Dimension & Candidate DOC measures \\
\midrule
Fidelity & compliance with constitutional, statutory, court, and department mandates \\
Consistency & comparable disciplinary, classification, grievance, and release decisions \\
Transparency & public reporting, records access, incident disclosure, and grievance visibility \\
Efficiency & case processing, staffing capacity, program delivery, and backlog measures \\
Equity & demographic disparity, disability access, language access, and geographic/facility differences \\
Accountability & corrective actions, inspections, investigations, remediation, and recurrence \\
\bottomrule
\end{tabular}
\end{center}

These are candidate families, not approved metrics. Each metric requires a
written definition, valid denominator, source inventory, normalization rule,
sufficiency threshold, and legal review before scoring.

\subsection{DOC-Specific Safeguards}

\begin{itemize}
  \item Separate department, facility, contractor, and individual officer
        entities.
  \item Distinguish allegations, grievances, investigations, findings, and
        adjudicated outcomes.
  \item Never use incarceration or demographic outcome differences without
        appropriate denominators and documented confounder handling.
  \item Treat missing or suppressed correctional data as a transparency and
        evidence-availability question, not automatic proof of misconduct.
  \item Preserve facility-level and population-level results separately.
\end{itemize}

The DOC profile is successful only if it runs through the same universal
entity, evidence, metric-observation, and score-result contracts as the NCRA
profile while producing a distinct, profile-versioned metric registry.
